\begin{enumerate}
    \item A l'aide d'une cellule, on détermine la conductance d'une solution
          $S_{1}$ de chlorure de sodium $\ce{NaCl}$ de concentration
          $c=\qty{5e-3}{\mol\per\L}$~; on trouve $G=\qty{5.45e-3}{\S}$.
          \begin{enumerate}
              \item Écrire la réaction de dissociation du chlorure de sodium
                    dans l'eau.
              \item La dissociation de $\ce{NaCl}$ est totale. Déterminer les
                    concentrations en $\unit{\mol\per\L}$, puis en
                    $\unit{\mol\per\cubic\m}$ des ions $\ce{Na+}$ et $\ce{Cl-}$
                    ($[\ce{Na+}]$ et $[\ce{Cl-}]$) dans la solution. La réponse
                    sera clairement justifiée.
              \item On donne les conductivités molaires ioniques l(Na+)= 3,87 10-3 .S.m2.mol-1 et l(Cl-)=7,63 10-3 S.m2.mol-1 . Déterminer la conductivité s de la solution .
              \item K= L/S (L~: distance entre les électrodes , S surface immergée d'une électrode) est appelée "constante de cellule" . Déterminer K .
          \end{enumerate}
    
    \item On dilue 10 fois la solution précédente (notée S1)~: on appelle S2 la solution obtenue .
          \begin{enumerate}
              \item Proposez un mode opératoire qui permette d'obtenir 100 mL de S2 à partir de la solution S1 .
              \item Quelles sont alors les concentrations des espèces ioniques présentes dans la solution S2~?
              
                  On utilise le même cellule conductimétrique que précedemment pour mesurer la conductance de la solution S2 .
              \item Déterminer la conductance de la solution S2
              \item La tension aux bornes de la cellule est égale à 1 V . Calculer l'intensité I du courant qui traverse la cellule sachant que la constante de cellule est la même que précedemment .
          \end{enumerate}
\end{enumerate}
